% SPECTRUM ICLR 2027 -- OPTIONAL REVIEW-STAGE APPENDIX, 2026-11-07.
% Do NOT replace the submitted fixed-tau method or the 1.5B five-round main table.
% Source status: audited local native evaluator summaries; raw files are not
% currently committed to GitHub. Verify original complete.json, model/config
% identity, task/sample denominator, and evaluator receipts before submission.
% This is a ready-to-integrate excerpt, not the current OpenReview PDF.
%
% Suggested ONE-LINE addition at the end of the RQ3 discussion in the main text:
%   A separate, post-submission 7B analysis of gain design and spectral strength
%   is reported in Appendix~\ref{app:postreview-gains}; its different model,
%   generation budget, and calibration access prevent pooling it with our
%   original 1.5B causal comparison.
%
% Append the following material AFTER \appendix and BEFORE \end{document}.

\section{Additional spectral-gain design and strength controls}
\label{app:postreview-gains}

\paragraph{Scope and comparability.}
The submitted primary method uses the fixed, proximal full-rank gain
$s_j=[1+\tau(1-\mu_j)]^{-1}$, with $\tau=1$. Here we report additional
native-student measurements recovered from a separately executed study
on Qwen2.5-Coder-7B-Instruct. These measurements are a
\emph{gain-design extension}, not a substitute for the submitted
1.5B five-round result or a proof that reference-gradient eigenvectors
are necessary. No evaluation outcome was used to select a generated
training record. Results below use the native student after five
updates and do not include generation-time spectral folding at test time.

\paragraph{Fixed proximal gains versus learned gains.}
The exploratory learned variant keeps the same loss-sensitive K/V
eigenspaces but optimizes continuous gains $a_j\in[0.1,1]$ on reference
completion likelihood. It validates a native-to-steered probe KL
budget of $0.02$ nats/token, checks a held-out reference guard,
and uses the accepted operator only for generation; the LoRA student
still learns from unfiltered raw completions. The learned variant uses
40 fit references, 10 independent guard references and 8 native
training-side probes. Thus it has a larger calibration/checking budget
than a fixed-gain implementation using 40 fit references, and
additional optimization compute. This comparison isolates a
\emph{family of gain choices} under a common backbone/benchmark,
not all other causes of difference.

\begin{table}[h]
\centering\small
\caption{\textbf{Exploratory five-round gain-design comparison
on Qwen2.5-Coder-7B-Instruct.} MBPP500; 16 training completions
per synthesis prompt per round; 4096 maximum new tokens,
6144 SFT tokens, 64 native evaluation draws per task.
Reported values are single-run point estimates, not a
statistical claim of method dominance. $C_{64}$ is the
mean number of distinct correct AST structures across
all 500 tasks. Learned has additional reference guard
and KL-probe cost. Correct-count-matched $D_4$ is
intentionally omitted until raw per-task records are
revalidated against its occupancy definition.}
\label{tab:postreview-gains}
\begin{tabular}{lrrr}
\toprule
Synthesis policy & pass@1 (\%) & pass@64 (\%) & $C_{64}$\\
\midrule
Plain SD                & \textbf{59.90} & 82.40 & 9.174\\
SSD decoding recipe     & 56.71 & 84.20 & 13.730\\
Fixed proximal SPECTRUM & 56.85 & \textbf{84.80} & \textbf{14.570}\\
Learned-gain extension  & 58.93 & 82.80 & 10.186\\
\bottomrule
\end{tabular}
\end{table}

Relative to fixed gains, learned gains improve pass@1 by
$2.07$ percentage points but reduce pass@64 by $2.0$ points
and $C_{64}$ by $4.384$ correct AST classes. In this matched
five-round setting, learning reference-likelihood-driven gains
moves the observed operating point toward single-draw accuracy,
rather than dominating fixed gains in correctness and structural
breadth. This finding is consistent with a nontrivial
accuracy--diversity trade-off, but it does \emph{not} establish
a causal benefit of the calibrated eigenbasis: the requisite
random-orientation and strength-matched isotropic training
comparisons remain outstanding. Nor is the single observed
trajectory evidence of across-training-seed stability.

\paragraph{Spectral-strength sensitivity.}
A separate \emph{single-round} Qwen2.5-Coder-7B study used a
512-token generation cap and 16 evaluation draws per MBPP task.
It must not be combined statistically with the five-round
4096-token table above.

\begin{table}[h]
\centering\small
\caption{\textbf{One-round fixed-gain strength sweep.}
Qwen2.5-Coder-7B, MBPP500, at most 512 generated tokens,
16 native evaluation draws per task. Values are unpaired
point estimates from separately collected sample pools;
no strength is selected as a new main result based on
test-set outcomes.}
\label{tab:postreview-tau}
\begin{tabular}{lrrrrrr}
\toprule
$\tau$ & 1.0 & 1.5 & 2.0 & 3.0 & 5.0 & 10.0\\
\midrule
pass@1 (\%) & 56.74 & 56.44 & 56.79 & 56.74 & 56.33 & 56.48\\
pass@16 (\%) & 76.60 & 76.40 & 77.20 & 76.60 & 75.60 & 77.40\\
\bottomrule
\end{tabular}
\end{table}

The single-round pass@16 estimates are not monotonic in $\tau$,
and higher $\tau$ does not consistently improve pass@1.
They demonstrate sensitivity to intervention strength, not
an optimal long-run choice; they do not isolate reference
information from geometric orientation. Separate ten-round
4096-token fixed-strength, decayed-strength, adaptive and
learned-gain experiments probe stability, but only the
learned and $\tau$-decay $0.8$ trajectories have complete
ten-round native receipts in the recovered audit.
Incomplete round-nine/ten cells remain unreported rather
than being inferred.

\paragraph{What remains unproven.}
Although one completed 1.5B single-round isotropic-gain
control exists in the archived results, its matched reference
is not an otherwise identical five-round fixed-$\tau=1$
trajectory. A completed random-eigenbasis training control,
anchor-only-information control, and complete multi-seed
fixed-gain confirmation are not established by these
supplementary results. Accordingly, we do not claim that
calibrated orientation is necessary, that learned gains
uniformly outperform the original fixed formulation, or
that differences between the two generation policies
come solely from gain adaptation.
